\chapter{容器}
\label{ch:11}

前面几章你已经偷偷用过容器了——猜数字里没有，但从第 5 章的 \texttt{List<\allowbreak{}int64>\allowbreak{}} 开始就一直在用。这一章把它讲清楚。

Magel 内建三种容器，不用 \texttt{import}：

\begin{center}
\begin{tabular}{lll}
\toprule
容器 & 装什么 & 什么时候用 \\
\midrule
\texttt{List<\allowbreak{}T>\allowbreak{}} & 有序的一串，按位置访问 & 顺序重要、要按下标取 \\
\texttt{Map<\allowbreak{}K,\allowbreak{} V>\allowbreak{}} & 键 \ensuremath{\rightarrow} 值 & 按名字查东西 \\
\texttt{Set<\allowbreak{}T>\allowbreak{}} & 不重复的一堆，自动按升序 & 去重、判断「有没有」 \\
\bottomrule
\end{tabular}
\end{center}

\section{\texttt{List<\allowbreak{}T>\allowbreak{}}：有序表}

\par
\medskip
\begin{lstlisting}[style=magel]
int64 main()
{
    List<int64> scores;
    scores.push(30);
    scores.push(10);
    scores.push(20);

    println("len={} at1={} has10={} idx={}",
        scores.len(), scores.at(1), scores.has(10), scores.indexOf(20));

    scores.insert(0, 7);
    println("insert 后：{} len={}", scores[0], scores.len());

    int64 last = scores.pop();
    println("pop={} len={}", last, scores.len());

    scores.remove(0);
    println("remove(0) 后：{} len={}", scores[0], scores.len());

    scores.clear();
    println("clear 后：len={}", scores.len());
    return 0;
}
\end{lstlisting}
\par
\medskip

\par
\medskip
\begin{lstlisting}[style=out]
len=3 at1=10 has10=1 idx=2
insert 后：7 len=4
pop=20 len=3
remove(0) 后：30 len=2
clear 后：len=0
\end{lstlisting}
\par
\medskip

常用操作：

\begin{center}
\begin{tabular}{lll}
\toprule
操作 & 作用 & 注意 \\
\midrule
\texttt{push(\allowbreak{}x)\allowbreak{}} & 在末尾加一个 & 最常用 \\
\texttt{insert(\allowbreak{}i,\allowbreak{} x)\allowbreak{}} & 在位置 \texttt{i} 插入 & 后面的元素整体后移 \\
\texttt{pop(\allowbreak{})\allowbreak{}} & 拿走末尾那个并返回 & 空表上调用会出错 \\
\texttt{remove(\allowbreak{}i)\allowbreak{}} & 删除位置 \texttt{i} & 按下标删，不是按值 \\
\texttt{clear(\allowbreak{})\allowbreak{}} & 清空 &  \\
\texttt{len(\allowbreak{})\allowbreak{}} & 元素个数 &  \\
\texttt{at(\allowbreak{}i)\allowbreak{}} / \texttt{list[\allowbreak{}i]\allowbreak{}} & 取位置 \texttt{i} & 两种写法等价 \\
\texttt{has(\allowbreak{}x)\allowbreak{}} & 有没有这个值 & 返回 \texttt{bool}（打印是 1/0） \\
\texttt{index\allowbreak{}Of(\allowbreak{}x)\allowbreak{}} & 第一次出现在哪 & 没有时返回什么见下 \\
\bottomrule
\end{tabular}
\end{center}

\texttt{index\allowbreak{}Of} 找不到时返回 \texttt{-1}——这是惯例，用的时候要判断：

\par
\medskip
\begin{lstlisting}[style=magel]
int64 position = scores.indexOf(99);
if (position >= 0)
{
    println("在位置 {}", position);
}
else
{
    println("没有这个值");
}
\end{lstlisting}
\par
\medskip

\subsection{越界会被抓到}

\par
\medskip
\begin{lstlisting}[style=magel]
int64 main()
{
    List<int64> numbers;
    numbers.push(1);
    println("{}", numbers[5]);
    return 0;
}
\end{lstlisting}
\par
\medskip

\par
\medskip
\begin{lstlisting}[style=out]
magel：运行期错误：vec 下标越界：5（长度 1）
出错位置：hello.mgl:4（main）
调用栈（外 → 内）：
  main  hello.mgl:4 [thread 0 coroutine 0 ip 25]
\end{lstlisting}
\par
\medskip

不是未定义行为，是明确的运行期错误，还告诉你「下标 5、长度 1」。比起读到一块随机的内存，这个错误至少能让你立刻定位。

\section{\texttt{Map<\allowbreak{}K,\allowbreak{} V>\allowbreak{}}：键 \ensuremath{\rightarrow} 值}

\par
\medskip
\begin{lstlisting}[style=magel]
int64 main()
{
    Map<String, int64> ages;
    ages.put("amy", 11);
    ages["cid"] = 33;

    println("len={} amy={} hasBob={} cid={}",
        ages.len(), ages["amy"], ages.has("bob"), ages["cid"]);

    println("查一个没有的键：{}", ages["zoe"]);

    for (int64 i = 0; i < ages.len(); i = i + 1)
    {
        println("{} -> {}", ages.keyAt(i), ages.valAt(i));
    }
    return 0;
}
\end{lstlisting}
\par
\medskip

\par
\medskip
\begin{lstlisting}[style=out]
len=2 amy=11 hasBob=0 cid=33
查一个没有的键：0
amy -> 11
cid -> 33
\end{lstlisting}
\par
\medskip

写法：

\begin{itemize}

\item \texttt{put(\allowbreak{}k,\allowbreak{} v)\allowbreak{}} 和 \texttt{m[\allowbreak{}k]\allowbreak{} =\allowbreak{} v} 是一回事，后者顺手一点；

\item \texttt{m[\allowbreak{}k]\allowbreak{}} 读；\texttt{has(\allowbreak{}k)\allowbreak{}} 判断有没有；\texttt{remove(\allowbreak{}k)\allowbreak{}} 删掉并返回是否删到了；

\item 遍历用 \texttt{key\allowbreak{}At(\allowbreak{}i)\allowbreak{}} / \texttt{val\allowbreak{}At(\allowbreak{}i)\allowbreak{}} 按下标走，\texttt{i} 从 0 到 \texttt{len(\allowbreak{})\allowbreak{} - 1}。Map 的遍历顺序是按升序键排列的，所以同一份数据每次遍历结果一致；

\end{itemize}

\begin{tip}{提示}
小心「查不存在的键」。 上面 \texttt{ages[\allowbreak{}"zoe"]\allowbreak{}} 返回了 \texttt{0}，而不是报错。因为 \texttt{V} 是 \texttt{int64}，缺省值就是 0。这不是「没有这个键」的意思，而是「我按 0 处理了」——想知道键在不在，先用 \texttt{has(\allowbreak{}k)\allowbreak{}} 问一次。这是 Map 最容易埋 bug 的地方。
\end{tip}

\section{\texttt{Set<\allowbreak{}T>\allowbreak{}}：不重复，自动升序}

\par
\medskip
\begin{lstlisting}[style=magel]
int64 main()
{
    Set<int64> unique;
    unique.add(5);
    unique.add(3);
    unique.add(5);
    // 重复的加不进去

    println("len={} has5={} has9={}", unique.len(), unique.has(5), unique.has(9));
    println("按升序看：at0={} at1={}", unique.at(0), unique.at(1));
    println("remove(5) 拿回 {}，现在 len={}", unique.remove(5), unique.len());
    return 0;
}
\end{lstlisting}
\par
\medskip

\par
\medskip
\begin{lstlisting}[style=out]
len=2 has5=1 has9=0
按升序看：at0=3 at1=5
remove(5) 拿回 1，现在 len=1
\end{lstlisting}
\par
\medskip

\texttt{Set} 的两个特点：

\begin{itemize}

\item 自动去重：\texttt{add(\allowbreak{}5)\allowbreak{}} 两次，长度还是 1；

\item 自动排序：遍历出来是升序（\texttt{at(\allowbreak{}0)\allowbreak{}} 是最小的）。

\end{itemize}

\texttt{remove} 返回的是「有没有删到」。所以「删一个可能不存在的元素」不用先查\texttt{has}，看返回值就知道。

\section{集合字面量}

不想一行行 \texttt{push} 的时候，可以直接写：

\par
\medskip
\begin{lstlisting}[style=magel]
int64 main()
{
    List<int64> order =
    {
        3, 1, 2
    };
    println("{} {}", order.len(), order[0]);

    Set<int64> numbers =
    {
        3, 1, 2
    };
    println("{}", numbers.len());

    Map<String, int64> table =
    {
        "a": 1, "b": 2
    };
    println("{} {}", table.len(), table["b"]);
    return 0;
}
\end{lstlisting}
\par
\medskip

\par
\medskip
\begin{lstlisting}[style=out]
3 3
3
2 2
\end{lstlisting}
\par
\medskip

同样一对花括号，左边的类型决定它是什么。

\begin{center}
\begin{tabular}{ll}
\toprule
左边写的类型 & \texttt{\{.\allowbreak{}.\allowbreak{}.\allowbreak{}\}} 的含义 \\
\midrule
\texttt{List<\allowbreak{}T>\allowbreak{}} & 有序表，元素按写的顺序排 \\
\texttt{Set<\allowbreak{}T>\allowbreak{}} & 去重的集合，自动升序 \\
\texttt{Map<\allowbreak{}K,\allowbreak{}V>\allowbreak{}} & 键值对（括号里带冒号） \\
\bottomrule
\end{tabular}
\end{center}

跨类型写会报错——\texttt{List<\allowbreak{}int64>\allowbreak{} v =\allowbreak{} \{"a":\allowbreak{} 1\};} 里带了冒号，那是 \texttt{Map} 的形态：

\par
\medskip
\begin{lstlisting}[style=out]
magel：集合字面量类型不匹配：map 不能初始化 List<int64>
\end{lstlisting}
\par
\medskip

空表也有讲究：\texttt{List<\allowbreak{}int64>\allowbreak{} v =\allowbreak{} \{\};} 是空表，\texttt{Set<\allowbreak{}int64>\allowbreak{} v =\allowbreak{} \{\};} 是空集，类型必须写清楚（\texttt{auto x =\allowbreak{} \{.\allowbreak{}.\allowbreak{}.\allowbreak{}\};} 这种不给上下文的老写法仍然当成 \texttt{Set}）。

要 \texttt{List} 也可以继续 \texttt{push}（或者用第 17 章的 \texttt{collect(\allowbreak{})\allowbreak{}}），两条路都通。

\section{数组：两种}

除了容器，还有两种数组：

\par
\medskip
\begin{lstlisting}[style=magel]
int64 main()
{
    int64 fixed[3];
    // 长度写在类型里
    fixed[0] = 1;
    fixed[1] = 2;
    fixed[2] = 3;
    println("{}", fixed[0] + fixed[1] + fixed[2]);

    int64 runtime[] = new[4];
    // 长度运行时决定
    runtime[0] = 10;
    println("{} {}", len(fixed), len(runtime));
    return 0;
}
\end{lstlisting}
\par
\medskip

\par
\medskip
\begin{lstlisting}[style=out]
6
3 4
\end{lstlisting}
\par
\medskip

\begin{center}
\begin{tabular}{lll}
\toprule
写法 & 长度 & 什么时候用 \\
\midrule
\texttt{T a[\allowbreak{}N]\allowbreak{}} & 编译期定死 & 大小你知道，要省内存、要对齐硬件 \\
\texttt{T a[\allowbreak{}]\allowbreak{} =\allowbreak{} new[\allowbreak{}n]\allowbreak{}} & 运行时决定 & 大小要算出来或者读进来 \\
\bottomrule
\end{tabular}
\end{center}

数组和 \texttt{List} 的区别：数组长度不可变，\texttt{List} 可以 \texttt{push} / \texttt{pop}。 日常业务逻辑几乎都用 \texttt{List}；数组留给「缓冲区」「硬件寄存器映射」「性能敏感的定长表」这类场合。

\texttt{len(\allowbreak{}a)\allowbreak{}} 对数组和容器都能用。越界同样会被检查到。

\section{别忘了：容器也是值}

第 6 章那条规则对容器照旧。这一点在容器上最容易踩：

\par
\medskip
\begin{lstlisting}[style=magel]
int64 main()
{
    List<int64> original;
    original.push(1);

    List<int64> copy = original;
    // 整表复制一份
    copy.push(2);

    println("original={} copy={}", original.len(), copy.len());
    return 0;
}
\end{lstlisting}
\par
\medskip

\par
\medskip
\begin{lstlisting}[style=out]
original=1 copy=2
\end{lstlisting}
\par
\medskip

\texttt{copy =\allowbreak{} original} 之后两张表毫无关系。 从 Java / Python 过来的人第一次写「把表传进函数想让它加元素」，会发现函数里加的东西外面看不到——

\par
\medskip
\begin{lstlisting}[style=magel]
void addOne(List<int64> numbers)
{
    numbers.push(1);
    // 改的是副本！
}
\end{lstlisting}
\par
\medskip

要真的改到外面那张表，写引用：

\par
\medskip
\begin{lstlisting}[style=magel]
void addOne(List<int64> mut& numbers)
{
    numbers.push(1);
}
// 调用：addOne(&numbers);
\end{lstlisting}
\par
\medskip

这也是性能上的关键点：容器一般都有一份堆上的数据，按值传一次就是整份复制。容器参数默认写 \texttt{const\&} 或 \texttt{mut\&}，除非你确实想要一份独立副本。

\subsection{容器里的东西也是副本}

\par
\medskip
\begin{lstlisting}[style=magel]
struct Player
{
    String name;
    int64 score;
};

int64 main()
{
    List<Player> players;
    Player p;
    p.name = "amy";
    p.score = 30;
    players.push(p);
    // 存进去的是一份副本

    p.score = 99;
    // 改原来的 p
    println("表里还是 {}", players[0].score);

    Player back = players[0];
    // 取出来的也是副本
    back.score = 1;
    println("表里还是 {}", players[0].score);
    return 0;
}
\end{lstlisting}
\par
\medskip

\par
\medskip
\begin{lstlisting}[style=out]
表里还是 30
表里还是 30
\end{lstlisting}
\par
\medskip

存进去是副本、取出来也是副本。 想原地改容器里的元素，得用引用（第 7 章）：\texttt{for (\allowbreak{}Player mut\& item in players)\allowbreak{}} 那类写法在讲完引用之后你就理解了。

\begin{tip}{提示}
字符串元素也一样：取出来的是新的一份，所以容器活得比某根字符串久也没关系。
\end{tip}

\subsection{容器归声明它的那一层}

容器在离开作用域时自动释放（第 8 章）：

\par
\medskip
\begin{lstlisting}[style=magel]
int64 main()
{
    {
        List<int64> inner;
        inner.push(1);
        inner.push(2);
        println("块里的表：len={}", inner.len());
    }
    // ← 这里把 inner 还回去了
    println("done");
    return 0;
}
\end{lstlisting}
\par
\medskip

\par
\medskip
\begin{lstlisting}[style=out]
块里的表：len=2
done
\end{lstlisting}
\par
\medskip

不用手动释放，也不用担心泄漏。要提前还就 \texttt{delete inner;}。

\section{遍历}

遍历用 \texttt{for-in}，不要数下标：

\par
\medskip
\begin{lstlisting}[style=magel]
int64 main()
{
    List<int64> scores;
    scores.push(12);
    scores.push(55);
    scores.push(41);

    int64 total = 0;
    int64 best = 0;
    for (int64 score in scores)
    {
        total = total + score;
        if (score > best)
        {
            best = score;
        }
    }
    println("总和 {} 最高 {}", total, best);
    return 0;
}
\end{lstlisting}
\par
\medskip

\par
\medskip
\begin{lstlisting}[style=out]
总和 108 最高 55
\end{lstlisting}
\par
\medskip

\texttt{for-in} 遍历出来的是元素的副本，所以「边遍历边改」改不到容器。要用下标改，就用 \texttt{for}：

\par
\medskip
\begin{lstlisting}[style=magel]
for (int64 i = 0; i < scores.len(); i = i + 1)
{
    scores[i] = scores[i] + 10;
    // 这样才改到表里
}
\end{lstlisting}
\par
\medskip

\begin{tip}{提示}
遍历时不要增删元素。 \texttt{for-in} 途中 \texttt{push} / \texttt{remove} 会让「下一个是谁」变得不可预测。要筛要删，用第 17 章的迭代器 \texttt{filter} / \texttt{collect} 造一张新表，或者先记下要删的下标再统一删。
\end{tip}

\subsection{用 \texttt{Any} 保存不同类型的元素}

\texttt{List<\allowbreak{}Any>\allowbreak{}} 明确要求元素装箱；\texttt{Set<\allowbreak{}Any>\allowbreak{}} 和 \texttt{Map} 的键或值槽也可以明确写成 \texttt{Any}。普通的 \texttt{List<\allowbreak{}int64>\allowbreak{}} 仍只接受能转换为 \texttt{int64} 的值。

\par
\medskip
\begin{lstlisting}[style=magel]
int64 main()
{
    List<Any> values = {1, "two", true};
    values.push(1.5);
    Any text = values[1];
    println("{} {}", values.len(), tryGet<String>(&text).unwrap());

    Set<Any> unique;
    Any key = "key";
    unique.add(key);
    unique.add(key);
    unique.add(2);
    println("{}", unique.len());

    Map<Any, Any> mixed;
    mixed.put(key, "stored");
    mixed.put(3, true);
    Any saved = mixed[key];
    println("{} {}", mixed.len(), tryGet<String>(&saved).unwrap());

    Map<String, Any> record = {"count": 3, "name": "Magel"};
    Any count = record["count"];
    println("{}", tryGet<int64>(&count).unwrap());
    saved.clear();
    println("{}", saved.isEmpty());
    return 0;
}
\end{lstlisting}
\par
\medskip

\par
\medskip
\begin{lstlisting}[style=out]
4 two
2
2 stored
3
1
\end{lstlisting}
\par
\medskip

\texttt{Any} 保留精确类型，\texttt{try\allowbreak{}Get<\allowbreak{}T>\allowbreak{}} 返回拥有值的副本；类型不匹配时返回 \texttt{None}。\texttt{clear(\allowbreak{})\allowbreak{}} 释放盒子当前载荷并恢复空状态，不影响之前复制出的值。字符串、嵌套容器按值复制，盒子和容器在离开作用域时自动清理。

\texttt{Any} 键比较身份，不比较装箱内容。 复制 \texttt{key} 保留身份，所以重复插入会去重，并且可以用它查询映射。两次独立装箱相同内容是不同键；需要按字符串内容查询时使用 \texttt{Map<\allowbreak{}String,\allowbreak{} Any>\allowbreak{}}。异构集合的排列也不要当成按载荷数值或文本排序。

装箱不会延长引用寿命。局部结构体、数组等多槽值的地址不能被保存在更长寿命的盒子里；需要接口视图时可直接绑定实现类型的引用，无须经过 \texttt{Any}。

\section{小结}

\begin{center}
\begin{tabular}{ll}
\toprule
容器 & 常用操作 \\
\midrule
\texttt{List<\allowbreak{}T>\allowbreak{}} & \texttt{push} \texttt{pop} \texttt{insert} \texttt{remove} \texttt{clear} \texttt{len} \texttt{at} \texttt{[\allowbreak{}]\allowbreak{}} \texttt{has} \texttt{index\allowbreak{}Of} \\
\texttt{Map<\allowbreak{}K,\allowbreak{}V>\allowbreak{}} & \texttt{put} \texttt{[\allowbreak{}]\allowbreak{}} \texttt{has} \texttt{remove} \texttt{len} \texttt{key\allowbreak{}At} \texttt{val\allowbreak{}At} \\
\texttt{Set<\allowbreak{}T>\allowbreak{}} & \texttt{add} \texttt{has} \texttt{remove} \texttt{len} \texttt{at} \\
\bottomrule
\end{tabular}
\end{center}

要记住的：

\begin{itemize}

\item 字面量：左边写 \texttt{List<\allowbreak{}T>\allowbreak{}} / \texttt{Set<\allowbreak{}T>\allowbreak{}} / \texttt{Map<\allowbreak{}K,\allowbreak{}V>\allowbreak{}} 决定 \texttt{\{.\allowbreak{}.\allowbreak{}.\allowbreak{}\}} 是什么，带冒号的是 \texttt{Map}；

\item 越界是明确的运行期错误，不是未定义行为；

\item \texttt{Map} 读不存在的键得到缺省值（不报错），先 \texttt{has} 再取；

\item 容器是值：赋值、传参、存取元素都是复制；传参默认写 \texttt{const\&} / \texttt{mut\&}；

\item 遍历用 \texttt{for-in}，遍历出来是副本；遍历时不要增删；

\item 容器在出作用域时自动释放。

\end{itemize}

\section{下一步}

现在你会用容器了，但要写「取分数前 3 名」「把所有名字变大写」这种活儿，用 \texttt{for-in} 加手工判断会很啰嗦。

不过在那之前，先解决一个你已经在用的东西但没细讲过的：字符串。它的下标是按字节的——这跟中文直接相关，值得单独一章。

翻到\hyperref[ch:12]{第 12 章 字符串与文本}。
